\documentclass[11pt,a4paper]{article}
\usepackage[top=0.75in,bottom=0.75in,left=0.75in,right=0.75in]{geometry}
\usepackage[T1]{fontenc}
\usepackage{times,microtype,enumitem,hyperref,titlesec}
\pagestyle{empty}
\setlength{\parindent}{0pt}
\setlength{\parskip}{0pt}
\linespread{1.10}
\hypersetup{hidelinks,breaklinks=true}
\titleformat{\section}{\bfseries\normalsize}{}{0pt}{}[\titlerule]
\titlespacing*{\section}{0pt}{10pt}{4pt}
\newcommand{\entry}[3]{\noindent\textbf{#1}\hfill\textbf{#3}\\[-1pt]\textit{#2}\par}
\newenvironment{cvitems}{\begin{itemize}[leftmargin=1.1em,label=\textbullet,itemsep=2pt,topsep=3pt,parsep=0pt,partopsep=0pt]}{\end{itemize}}
\begin{document}
{\LARGE\bfseries Aarush Ram Anandh}\par\vspace{4pt}
\href{mailto:aarush.anandh_ug2024@ashoka.edu.in}{aarush.anandh\_ug2024@ashoka.edu.in}\quad|\quad +91 9659786783\par\vspace{2pt}
\href{https://github.com/MushuRam01}{github.com/MushuRam01}\quad|\quad
\href{https://www.linkedin.com/in/aarush-ram-anandh/}{linkedin.com/in/aarush-ram-anandh}\par
\vspace{20pt}
11 November 2026\par\vspace{12pt}
Hiring Team\\
GE HealthCare\par\vspace{12pt}
\textbf{Application: Research Intern -- AI}\par\vspace{14pt}
Dear Hiring Team,\par\vspace{10pt}
I am applying for the Research Intern -- AI role at GE HealthCare. I am completing a B.Sc. (Hons.) in Computer Science and Mathematics at Ashoka University, expected in 2028. The role\textquotesingle s combination of image analytics, generative models and medical-workflow research connects with my experience in diffusion models, computer vision and physiological time-series analysis.\par\vspace{10pt}
At the Constrained Image-Synthesis Lab, I worked on colour preservation in latent diffusion, achieving $\Delta E_{00}<1$ and approximately four times better colour accuracy than evaluated open-source baselines. The methods and results are described in my public preprint, \textit{Inference-Time Loss-Guided Colour Preservation in Diffusion Sampling} (arXiv:2601.17259). My ongoing CGSW research investigates controllable optimal-transport distances, with a paper submitted to JMLR and a hyperspectral experiment on 224-band images across seven waste categories. These projects have given me experience turning a mathematical objective into an implementation and an evaluation.\par\vspace{10pt}
At the Koita Centre for Digital Health, I contributed to a wearable-data pipeline achieving approximately 70\% sensitivity (recall) for asthma exacerbation detection in study evaluation. I developed patient-day processing for 191 patients and 15,691 patient-days, with HR/HRV smoothing, residuals, change-point features and tree-based and recurrent modelling workflows. The evaluation separates same-day detection from future-onset prediction and uses patient-wise splits across five seeds. At Avartan Labs, I raised rare-class correct detection from 40\% to 80\% in an 18-class RT-DETR detection and tracking pipeline. These experiences connect physiological-data research with practical image analytics.\par\vspace{10pt}
I would welcome the opportunity to contribute Python, C++ and PyTorch implementation skills to research on image quality and AI-assisted workflows, while learning from your technologists and engineers. I would bring experience with noisy data, model failure analysis and research prototyping to the team. Thank you for considering my application.\par\vspace{14pt}
Yours sincerely,\\[5pt]
Aarush Ram Anandh
\end{document}
